% LPC 2026, Tracing Microconference: three slides. The PDF goes to the
% organizers, is published, and is what the room's projector shows: the slide
% on upstream involvement is therefore in the deck, not held back on a laptop.
% The backup slide on other kernel policies is parked in backup.tex.
%
% The content is plan.txt, in this directory. The look (fonts, palette, kicker
% above the title, footline) is ../lpc2026-slides', copied rather than shared:
% neither deck can then break the other.

\documentclass[aspectratio=169, 10pt, t]{beamer}

% ---------------------------------------------------------------------------
% Fonts. IBM Plex is the design's face; where it is not installed, Lato + Noto
% Sans Mono stand in. A font swap can move a line break, so `make` reports any
% overfull box.
% ---------------------------------------------------------------------------
\usepackage{fontspec}
\IfFontExistsTF{IBM Plex Sans}{%
  % The 600 weight's name depends on the build: upstream's OTFs call it
  % "SemiBold", Debian's fonts-ibm-plex TTFs "SmBld".
  \IfFontExistsTF{IBM Plex Sans SemiBold}{\def\plexsb{SemiBold}}{\def\plexsb{SmBld}}%
  \setsansfont{IBM Plex Sans}[BoldFont={IBM Plex Sans \plexsb},
                              BoldItalicFont={IBM Plex Sans \plexsb\space Italic}]
  \setmonofont{IBM Plex Mono}[BoldFont={IBM Plex Mono \plexsb}]
}{%
  \setsansfont{Lato}[BoldFont={Lato Semibold},
                     BoldItalicFont={Lato Semibold Italic}]
  \setmonofont{Noto Sans Mono}[Scale=0.945,
                               BoldFont={Noto Sans Mono Bold}]
}

\usepackage{tikz}
\usepackage{tcolorbox}
\tcbuselibrary{skins}

% ---------------------------------------------------------------------------
% Palette.
% ---------------------------------------------------------------------------
\definecolor{paper}{HTML}{F5F3EE}
\definecolor{paperalt}{HTML}{E9E5DC}
\definecolor{ink}{HTML}{1A1D23}
\definecolor{muted}{HTML}{4F5561}
\definecolor{accent}{HTML}{2456A6}

% ---------------------------------------------------------------------------
% Beamer look: quiet. A small section kicker above each title stands in for
% divider slides.
% ---------------------------------------------------------------------------
\setbeamersize{text margin left=9mm, text margin right=9mm}
\setbeamertemplate{navigation symbols}{}
\setbeamercolor{background canvas}{bg=paper}
\setbeamercolor{normal text}{fg=ink, bg=paper}
\setbeamercolor{structure}{fg=accent}
\setbeamercolor{alerted text}{fg=accent}
\setbeamercolor{frametitle}{fg=ink}
\setbeamercolor{kicker}{fg=muted}
\setbeamercolor{footline}{fg=muted}
\setbeamerfont{frametitle}{size=\LARGE, series=\bfseries}
\setbeamerfont{kicker}{size=\scriptsize, series=\bfseries}
\setbeamerfont{footline}{size=\tiny}

% \talksection sets both beamer's section and the kicker. \gdef, and a plain
% \def for the empty case: the template tests it against \@empty with \ifx,
% which a \newcommand-defined (\long) macro would never equal.
\def\kickertext{}
\newcommand{\talksection}[1]{\section{#1}\gdef\kickertext{#1}}

\makeatletter
\setbeamertemplate{frametitle}{%
  \vskip 4mm
  \ifx\kickertext\@empty\else
    {\usebeamerfont{kicker}\usebeamercolor[fg]{kicker}%
     \addfontfeature{LetterSpace=6}\MakeUppercase{\kickertext}\par}%
    \vskip 1mm
  \fi
  {\usebeamerfont{frametitle}\usebeamercolor[fg]{frametitle}%
   \insertframetitle\par}%
  \vskip 1mm}
\makeatother

% The EfficiOS wordmark, for the footline. efficios-logo.pdf is vector art on
% a letter page, and the viewport boxes it: left and right at the ink's edges,
% top at the top of the E, bottom at the letters' BASELINE (y = 424.1 bp; the
% ff descends to 405.6). It is not clipped, so the descenders hang under the
% box as a glyph's would: the E starts on the text margin and the wordmark
% shares the footline's baseline. height= is therefore the height of the E.
% \smash: the footline stays as tall as its text, so the logo moves no slide's
% layout.
\newcommand{\footlogo}{\smash{%
  \includegraphics[viewport=149.4 424.1 457.3 501.2, height=3.7mm]%
    {efficios-logo.pdf}}}

% No cover slide in a microconference slot, so the footline names the speaker;
% the logo stands for the affiliation.
\setbeamertemplate{footline}{%
  \usebeamerfont{footline}\usebeamercolor[fg]{footline}%
  \hbox to \paperwidth{\hskip 9mm \footlogo\hskip 2.5mm
    Mathieu Desnoyers\enspace·\enspace
    Tracing Microconference\enspace·\enspace LPC 2026\hfill
    \insertframenumber\hskip 9mm}%
  \vskip 3.5mm}

\setbeamertemplate{itemize item}{\textcolor{accent}{–}}
\setbeamertemplate{itemize subitem}{\textcolor{muted}{–}}
\setlength{\leftmargini}{1.1em}

% ---------------------------------------------------------------------------
% Text blocks.
% ---------------------------------------------------------------------------
\tcbset{calloutbase/.style={enhanced, colback=paperalt, frame hidden, arc=2pt,
  borderline west={1.6pt}{0pt}{accent}, left=3mm, right=3mm, top=1.5mm,
  bottom=1.5mm, before skip=1mm, after skip=1mm, halign=flush left}}
\newtcolorbox{callout}[1][]{calloutbase, fontupper=\small, #1}

% A numbered constraint: the number hangs in the left margin of its box.
\newlength{\constraintindent}
\setlength{\constraintindent}{11mm}
\newtcolorbox{constraint}[1]{calloutbase,
  fontupper=\fontsize{11}{14}\selectfont,
  enlarge left by=\constraintindent, width=\linewidth-\constraintindent,
  overlay={\node[anchor=north west, inner sep=0pt, text=accent,
                 font=\fontsize{26}{26}\selectfont\bfseries]
           at ([xshift=-\constraintindent, yshift=-0.6mm]frame.north west) {#1};}}
% Who said it, under a constraint.
\newcommand{\source}[1]{{\par\vspace{1mm}\footnotesize\color{muted}#1\par}}

% A bold heading over a short list.
\newcommand{\area}[1]{{\par\vspace{1.3mm}\textbf{#1}\par}}

\title{LPC 2026 Tracing Microconference}
\author{Mathieu Desnoyers}
\institute{EfficiOS Inc.}
\date{Linux Plumbers Conference 2026}

\begin{document}

% ===========================================================================
% 1. The two constraints
% ===========================================================================
\talksection{The constraints}

\begin{frame}{Two constraints on kernel tracers}
  \begin{constraint}{1}
    Tracers are core code: they should be \textbf{built-in code only}, they
    have no justification for being kernel modules.
    \source{Paraphrasing core kernel developers}
  \end{constraint}
  \vspace{1mm}
  \begin{constraint}{2}
    ``If people need two tracing models, then the other one will be out of
    tree. It's that simple.\par
    \vspace{1.5mm}
    So as far as I'm concerned, this discussion is not a discussion. Either
    there's a way to merge things incrementally with SHARED infrastructure,
    or there isn't.''
    \source{Linus Torvalds, LKML, July 2025\\[0.3mm]
      {\tiny\url{https://lore.kernel.org/lkml/CAHk-=wiT9Cz+EbbuKozqiu7DnZQ7ftAWSmGf-xy_CdhJPCsNSg@mail.gmail.com/}}}
  \end{constraint}
\end{frame}

% ===========================================================================
% 2. What the two do together
% ===========================================================================
\talksection{The consequence}

\begin{frame}{Combined: a collaboration inhibitor}
  \fontsize{11}{14}\selectfont
  \begin{itemize}\setlength{\itemsep}{3mm}
  \item It is more efficient to work \textbf{outside of common tracer code}
        if that code cannot be used by an out-of-tree kernel tracer.
  \item It actively \textbf{prevents} incrementally contributing LTTng
        facilities into common tracer code \textbf{while keeping those
        facilities usable by the LTTng tracer}.
  \item Making perf, ftrace or eBPF use those facilities may justify
        integrating those contributions, but \textbf{does not justify
        exporting the required symbols} to GPL~modules.
  \end{itemize}
  \vfill
  \begin{callout}[fontupper=\large\bfseries, top=2mm, bottom=2mm]
    What is doable is an \textcolor{accent}{all-or-nothing} tracer ABI
    swap,\\ \textcolor{accent}{never incremental}.
  \end{callout}
  \vspace{2mm}
\end{frame}

% ===========================================================================
% 3. Upstream involvement
% ===========================================================================
\talksection{Upstream involvement}

\begin{frame}{20 years of Linux kernel contributions}
\begin{columns}[T, onlytextwidth]
\column{0.54\textwidth}
  \area{Atomic and local atomic headers}
  \begin{itemize}
  \item Generalization (\textasciitilde 2007)
  \end{itemize}
  \area{asm goto / static keys}
  \begin{itemize}
  \item Active involvement in the original desiderata leading to GCC support
        of asm goto
  \end{itemize}
  \area{Tracing subsystem}
  \begin{itemize}
  \item Tracepoint implementation
  \item Faultable system call tracepoints
  \item Participation in the SFrame kernel vs userspace integration effort
        (stack walk)
  \item Reviewer
  \end{itemize}
\column{0.41\textwidth}
  \area{RCU subsystem}
  \begin{itemize}
  \item Membarrier
  \item Hazard pointers (ongoing)
  \item Reviewer
  \end{itemize}
  \area{Scheduler}
  \begin{itemize}
  \item Restartable Sequences
  \item Concurrency IDs
  \end{itemize}
  \area{Memory management}
  \begin{itemize}
  \item Hierarchical per-CPU RSS accounting (ongoing)
  \end{itemize}
\end{columns}
\end{frame}

\end{document}
